\chapter{Expérience 3 : Ingénierie de l'évaluation - La « Notation à l'Enveloppe » (Méthode Italienne)}

\section{Axe 1 : Intentions pédagogiques et fondements éthiques}

\subsection{Genèse et conception originale : la « Notation à l'Enveloppe »}

Le dispositif d'évaluation appelé \textbf{« Notation à l'Enveloppe »} (surnommé informellement « Méthode Italienne » en référence à une expérience personnelle de Gentil Organisateur en hôtel-club en Sicile, où l'équipe recevait une enveloppe globale de pourboires et gratifications à répartir de manière concertée entre animateurs selon l'engagement de chacun) est une ingénierie d'évaluation originale que j'ai transposée dans le champ pédagogique pour répondre aux écueils de la notation de groupe indifférenciée.

Ayant développé ce protocole de manière empirique pour résoudre la problématique d'équité identifiée lors des projets collaboratifs (notamment à la suite de l'Expérience 2), j'ai ensuite entrepris une recherche approfondie dans la littérature didactique pour en analyser la portée. J'ai alors constaté avec un grand intérêt que des équipes universitaires et chercheurs avaient expérimenté des démarches d'évaluation par les pairs fortement convergentes :
\begin{itemize}[leftmargin=*]
    \item \textbf{À l'UTC Compiègne et à l'Université de Poitiers} \parencite{evaluation_pairs_utc,evaluation_pairs_univ_poitiers}, des enseignants-chercheurs ont conçu des systèmes d'allocation de points intra-groupe pour responsabiliser les étudiants et mesurer la contribution réelle au projet.
    \item \textbf{Dans la recherche internationale sur le Peer Assessment} \parencite{topping_peer_assessment,scallon_evaluation}, des travaux théorisent le rôle de l'évaluation par les pairs dans la construction de l'auto-régulation et du sentiment de justice redistributive.
\end{itemize}

Cette convergence entre mon intuition pédagogique de terrain et les travaux de recherche universitaires confirme la pertinence et la transférabilité de cette ingénierie d'évaluation régulatrice.

\subsection{Problématique d'ingénierie d'évaluation et double constat}

Le développement de cette méthode répond à deux défis majeurs en ingénierie de formation :
\begin{enumerate}[leftmargin=*]
    \item \textbf{Le risque d'injustice redistributive} : L'attribution d'une note unique globale occulte les écarts d'investissement individuel, générant de la frustration chez les apprenants moteurs et un risque de déresponsabilisation.
    \item \textbf{Les limites de l'observation magistrale directe} : Dans des dispositifs de pédagogie active (projets en sous-groupes), le formateur ou enseignant ne peut appréhender seul la dynamique fine des interactions et la répartition effective des charges de travail.
\end{enumerate}

Cette expérimentation s'inscrit donc au cœur des compétences du Master MEEF PIF / FFMS : l'\textbf{ingénierie des dispositifs d'évaluation régulatrice} \parencite{perrenoud_pratique_reflexive,nicol_formative_assessment}, visant à transformer l'évaluation sommative subie en un vecteur d'auto-régulation, de métacognition et d'apprentissage sociocognitif.

\subsection{Le prérequis : maturité du groupe-classe}

Ce système n'est pas applicable à n'importe quelle classe. Il nécessite deux conditions préalables \parencite{evaluation_pairs_utc} :

\begin{itemize}[leftmargin=*]
    \item \textbf{Connaissance mutuelle} : Les élèves doivent se connaître suffisamment pour évaluer objectivement le travail de leurs pairs
    \item \textbf{Maturité collective} : Il doit exister une certaine confiance et une culture de l'honnêteté au sein du groupe-classe pour que le processus soit sain et constructif
\end{itemize}

Pour ces raisons, je ne mets en place cette méthode qu'à partir du \textbf{second semestre de première NSI}, après plusieurs mois de travail collectif ayant permis de développer cette maturité.

\subsection{Les objectifs pédagogiques implicites}

Au-delà de l'évaluation sommative, cette méthode vise plusieurs objectifs formatifs essentiels \parencite{perrenoud_pratique_reflexive} :

\subsubsection{Développer la métacognition}

Amener les élèves à réfléchir sur :
\begin{itemize}[leftmargin=*]
    \item Leur propre contribution au projet (auto-évaluation)
    \item La contribution des autres membres (évaluation par les pairs) \parencite{evaluation_pairs_univ_poitiers}
    \item Les critères d'une contribution de qualité (construction d'une grille d'évaluation mentale)
\end{itemize}

Cette dimension métacognitive est reconnue comme essentielle pour le développement de l'autonomie et de la pratique réflexive \parencite{perrenoud_pratique_reflexive,schon_reflective_practitioner}, et elle s'appuie sur les principes de l'évaluation formative et de l'auto-régulation des apprentissages \parencite{nicol_formative_assessment,topping_peer_assessment}.

\subsubsection{Enseigner la responsabilité}

Chaque membre est responsable non seulement de sa propre note, mais aussi de l'impact de son travail (ou de son absence de travail) sur le groupe. Cette responsabilisation développe une conscience professionnelle transférable au monde du travail \parencite{competences_transversales}.

\subsubsection{Former au management équitable}

Les chefs de projet sont initiés à la difficile tâche d'évaluer le travail de leurs pairs de manière juste et argumentée. Cette compétence managérial est rarement enseignée au lycée mais constitue un atout majeur pour leur future carrière \parencite{knowles_andragogy}.

\section{Axe 2 : Déroulement méthodologique en 5 étapes structurées}

La méthode de « Notation à l'Enveloppe » suit un protocole rigoureux en 5 étapes, garantissant à la fois l'équité du processus et l'acceptation du résultat par les élèves.

\subsection{Étape 1 : Présentation transparente des règles du jeu}

Lors de la présentation du premier projet de groupe de l'année, j'explique clairement et précisément le système d'évaluation \parencite{evaluation_pairs_utc} :

\begin{itemize}[leftmargin=*]
    \item \textbf{Principe fondamental} : La note de groupe sera distribuée entre les membres selon leur contribution effective
    \item \textbf{Rôle du chef de projet} : Il sera responsable de proposer la répartition des points, non pas en tant que "juge" mais en tant que "témoin privilégié" du travail de chacun
    \item \textbf{Mécanisme de validation} : Toutes les propositions seront discutées collectivement entre chefs de projet sous ma supervision
    \item \textbf{Objectifs pédagogiques} : Développer la responsabilité individuelle, la justice dans le travail collaboratif et la capacité d'auto-évaluation
\end{itemize}

Cette transparence initiale génère une compréhension et une adhésion des élèves, condition sine qua non du succès de la méthode.

\subsection{Étape 2 : Attribution d'un projet de groupe structuré}

Le projet assigné doit présenter certaines caractéristiques pour que la méthode soit pertinente :

\begin{itemize}[leftmargin=*]
    \item \textbf{Durée suffisante} : Au minimum 3 semaines pour permettre l'observation des contributions sur la durée
    \item \textbf{Tâches diversifiées} : Mélange de recherche, conception, réalisation technique, présentation orale
    \item \textbf{Interdépendance des rôles} : Chaque membre doit avoir une contribution identifiable mais complémentaire
    \item \textbf{Groupes de 3 à 4 élèves maximum} : Au-delà, l'évaluation par les pairs devient trop complexe
\end{itemize}

Observation significative : lors de la désignation des chefs de projet, les groupes choisissent leurs chefs en fonction de leur fiabilité plutôt que de leur popularité. Cette dynamique montre que le système répond à un besoin de justice ressenti par les élèves eux-mêmes \parencite{evaluation_pairs_univ_poitiers}.

\subsection{Étape 3 : Calcul et attribution du "capital points" à distribuer}

À l'issue du projet, j'évalue le travail collectif du groupe et attribue une note globale (par exemple : 12/20). Cette note génère un \textbf{"capital points"} que le chef de projet devra répartir.

\subsubsection{Exemple concret de calcul}

Pour un groupe de 3 élèves ayant obtenu 12/20 :
\begin{itemize}[leftmargin=*]
    \item \textbf{Capital total} : 12 × 3 = 36 points à distribuer
    \item \textbf{Contrainte} : La somme des notes individuelles doit égaler exactement 36 points
    \item \textbf{Amplitude possible} : Les notes peuvent varier de 8/20 à 16/20 environ (selon les cas)
\end{itemize}

Cette contrainte mathématique simple garantit que l'évaluation reste dans des proportions raisonnables tout en permettant une réelle différenciation.

\subsection{Étape 4 : Répartition interne par le chef de projet}

Le chef de projet dispose d'une \textbf{grille d'évaluation simplifiée} pour objectiver sa proposition \parencite{evaluation_formative}. Cette grille se concentre sur quatre critères principaux :

\begin{enumerate}[leftmargin=*]
    \item \textbf{Assiduité et ponctualité} : Présence effective aux séances de travail en groupe
    \item \textbf{Qualité des contributions} : Pertinence et finition du travail produit
    \item \textbf{Respect des délais} : Livraison des tâches assignées dans les temps convenus
    \item \textbf{Esprit collaboratif} : Entraide, communication constructive, gestion des conflits
\end{enumerate}

\subsubsection{Processus de réflexion guidé}

Le chef de projet doit :
\begin{itemize}[leftmargin=*]
    \item \textbf{Auto-évaluer} sa propre contribution en premier
    \item \textbf{Justifier} chaque attribution par des faits observables
    \item \textbf{Proposer} une répartition équitable en points entiers
    \item \textbf{Anticiper} les possibles contestations et préparer ses arguments
\end{itemize}

Cette étape développe des compétences managériales et d'analyse objective cruciales pour la suite de leur parcours.

\subsection{Étape 5 : "Table Ronde" de validation collective (Analyse des Postures)}

La réunion finale constitue le moment pédagogique le plus riche. Analysée à travers le modèle de \textbf{l'ajustement pédagogique} d'Éric Saillot et des postures de Bucheton \parencite{saillot_ajustement,bucheton_postures}, cette phase illustre la complexité de l'étayage enseignant.

Je ne suis pas simplement un "modérateur", je construis une \textbf{posture d'ajustement} systémique articulant le "penser-dire-faire", en naviguant entre trois postures clés :

\begin{itemize}[leftmargin=*]
    \item \textbf{Posture d'accompagnement (Aide latérale)} : C'est ma posture dominante. Je suis assis \textit{avec} les élèves, pas \textit{face} à eux. J'écoute, je reformule leurs arguments ("Si je comprends bien, tu valorises son effort plus que son résultat ?"), aidant à expliciter leur pensée sans imposer la mienne.

    \item \textbf{Posture de contrôle (Cadrage serré)} : J'active cette posture instantanément si l'éthique est menacée (ex: règlement de compte personnel, discrimination). Je deviens alors le garant des règles du jeu et de la sécurité psychologique de chacun ("Stop, cet argument n'est pas recevable car il ne porte pas sur le travail").

    \item \textbf{Posture de lâcher-prise (Autonomie confiée)} : À certains moments, je me mets volontairement en retrait ("Et vous, les autres chefs, qu'en pensez-vous ?") pour laisser la régulation sociale opérer. C'est le moment où la compétence s'incarne véritablement chez les élèves.
\end{itemize}

Cette flexibilité posturale permet de transformer une simple réunion administrative en un véritable espace de développement professionnel pour les élèves.

\section{Axe 3 : Résultats observés et impacts}

\subsection{Analyse du Multi-agenda (Bucheton) sur l'ensemble de la séquence}

Pour compléter cette analyse, je peux lire l'ensemble de ce dispositif à travers les cinq préoccupations du \textbf{multi-agenda} \parencite{bucheton_postures} :
\begin{itemize}[leftmargin=*]
    \item \textbf{Pilotage} : Le dispositif est très structuré (étapes, grilles, capital points) pour éviter la dérive "subjective".
    \item \textbf{Atmosphère} : L'enjeu est de maintenir un climat de confiance et de justice ("ethos").
    \item \textbf{Tissage} : Je relie constamment leur expérience immédiate (noter un camarade) à des enjeux professionnels (évaluation annuelle, travail d'équipe).
    \item \textbf{Étayage} : Je fournis les outils (la grille) pour qu'ils puissent réussir la tâche complexe d'évaluation.
    \item \textbf{Objets de savoir} : L'objet n'est plus le code informatique, mais la compétence métacognitive d'évaluation elle-même.
\end{itemize}

\subsection{Une différenciation modérée et acceptée}

\subsubsection{Données quantitatives}

Sur 5 ans d'utilisation de cette méthode, j'observe des écarts modérés et acceptés \parencite{evaluation_pairs_univ_poitiers}, révélant l'équité du processus.

\subsubsection{Acceptation qualitative}

La majorité des élèves \textbf{comprennent et acceptent leur note}, même lorsqu'elle est inférieure à celle du groupe. Cette acceptation s'explique par la transparence du processus et la légitimité du chef de projet \parencite{evaluation_formative}.

\subsection{Un leadership basé sur la fiabilité plutôt que sur la popularité}

Les élèves comprennent rapidement que choisir un chef « copain » mais peu fiable conduit à un résultat médiocre. Ils développent donc une approche plus pragmatique du choix de leur chef, valorisant la \textbf{fiabilité} et la \textbf{capacité d'organisation} \parencite{competences_transversales}.

\subsection{Impact puissant sur la motivation et l'inclusion}

Le système agit comme un outil de diagnostic et de motivation \parencite{pedagogie_active}, permettant à chaque élève de se sentir reconnu pour sa contribution.



\subsection{Matériels produits}

Pour soutenir cette démarche, j'ai conçu et mis à jour la fiche d'évaluation distribuée aux apprenants (\textbf{document original joint au format PDF :} \url{sources/grille.pdf}, également reproduit dans les Annexes). Ces outils constituent une ingénierie d'évaluation opérationnelle, affinée sur le terrain et transférable à d'autres dispositifs de formation par projet.

\subsubsection{Principe de la notation à l'enveloppe}
L'évaluation repose sur le principe de la "notation à l'italienne" :
\begin{itemize}
    \item L'enseignant attribue une \textbf{enveloppe globale de points} au groupe, basée sur sa grille d'évaluation du projet.
    \item Les étudiants (pairs) proposent une \textbf{répartition de ces points} entre eux, en fonction de l'implication individuelle de chacun.
\end{itemize}

Il n'y a donc pas de pondération fixe, mais une distribution dynamique des points alloués au collectif.

\subsubsection{Grille d'évaluation professorale}
Cette grille me permet d'évaluer le produit fini rendu par le groupe et de déterminer l'enveloppe de points à distribuer.

\begin{longtable}{|p{4cm}|p{2cm}|p{2cm}|p{2cm}|p{2cm}|p{2cm}|}
\caption{Grille d'évaluation "à l'italienne" - Évaluation professorale} \\
\hline
\textbf{Critères projet} & \textbf{Insuffisant (0-4)} & \textbf{Fragile (5-8)} & \textbf{Satisfaisant (9-12)} & \textbf{Très bien (13-16)} & \textbf{Excellent (17-20)} \\
\hline
\endfirsthead
\hline
\textbf{Critères projet} & \textbf{Insuffisant (0-4)} & \textbf{Fragile (5-8)} & \textbf{Satisfaisant (9-12)} & \textbf{Très bien (13-16)} & \textbf{Excellent (17-20)} \\
\hline
\endhead
\textbf{Qualité technique} &
Produit non fonctionnel &
Fonctionnalités basiques &
Produit fonctionnel stable &
Qualité technique élevée &
Excellence technique \\
\hline
\textbf{Innovation et créativité} &
Aucune originalité &
Quelques éléments originaux &
Approche créative &
Solutions innovantes &
Innovation remarquable \\
\hline
\textbf{Respect des contraintes} &
Contraintes ignorées &
Respect partiel &
Contraintes respectées &
Dépassement des attentes &
Excellence dans les contraintes \\
\hline
\textbf{Documentation} &
Aucune documentation &
Documentation minimale &
Documentation correcte &
Documentation complète &
Documentation exemplaire \\
\hline
\end{longtable}

\subsubsection{Grille d'évaluation par les pairs}
Cette grille est utilisée par le chef de projet pour évaluer l'implication de chaque membre.

\begin{longtable}{|p{4cm}|p{3cm}|p{3cm}|p{3cm}|p{3cm}|}
\caption{Grille d'évaluation "à l'italienne" - Évaluation par les pairs} \\
\hline
\textbf{Critères implication} & \textbf{Passif (0-5)} & \textbf{Ponctuel (6-10)} & \textbf{Actif (11-15)} & \textbf{Leader (16-20)} \\
\hline
\endfirsthead
\hline
\textbf{Critères implication} & \textbf{Passif (0-5)} & \textbf{Ponctuel (6-10)} & \textbf{Actif (11-15)} & \textbf{Leader (16-20)} \\
\hline
\endhead
\textbf{Participation aux réunions} &
Absent ou silencieux &
Présent mais peu participatif &
Participation régulière &
Animation et leadership \\
\hline
\textbf{Contribution technique} &
Aucune contribution &
Contributions mineures &
Contributions significatives &
Contributions clés et expertises \\
\hline
\textbf{Esprit d'équipe} &
Individualiste, conflits &
Coopération minimale &
Bon esprit d'équipe &
Moteur de l'équipe \\
\hline
\textbf{Respect des délais} &
Retards systématiques &
Quelques retards &
Ponctualité respectée &
Anticipation et organisation \\
\hline
\textbf{Aide aux autres} &
Aucune entraide &
Aide ponctuelle &
Entraide régulière &
Mentoring et formation \\
\hline
\end{longtable}

\subsubsection{Avantages pédagogiques}
Cette méthode structurée offre de multiples avantages :
\begin{itemize}
    \item Développement de la métacognition par l'auto-évaluation et l'évaluation des autres.
    \item Responsabilisation accrue des apprenants.
    \item Équité dans la reconnaissance des contributions individuelles.
    \item Préservation de la motivation collective.
    \item Formation à une évaluation objective et constructive.
\end{itemize}



\section{Axe 4 : Analyse réflexive approfondie et transférabilité}

\subsection{Analyse réflexive approfondie}

\subsubsection{Analyse a priori : Le pari de la confiance}
En concevant ce dispositif de "notation à l'enveloppe", je faisais le pari risqué de la \textbf{maturité des élèves}. Mon analyse a priori reposait sur l'hypothèse que :
\begin{itemize}
    \item Les élèves sont capables d'objectivité si on leur fournit des critères clairs (la grille).
    \item Le sentiment de justice est un moteur plus puissant que la camaraderie complaisante.
    \item La responsabilisation (être chef de projet) induit un changement de posture immédiat.
\end{itemize}

\subsubsection{Analyse a posteriori : La justice perçue}
L'expérience a globalement validé ces hypothèses, mais avec une nuance importante : la \textbf{justice doit être visible}. J'ai observé que les tensions ne naissaient pas de la note elle-même, mais du sentiment d'arbitraire. Dès lors que le chef de projet pouvait justifier sa décision par des faits notés dans la grille, la tension retombait. L'apprentissage majeur pour moi a été de comprendre que l'évaluation n'est pas seulement une mesure, c'est un acte de communication.

\subsubsection{Difficultés imprévues et régulations}
La mise en œuvre a révélé des frictions humaines que la théorie n'avait pas anticipées :

\begin{itemize}[leftmargin=*]
    \item \textbf{Anxiété du "juge"} : Certains chefs de projet étaient paralysés par la peur de briser des amitiés.
    \item \textbf{Solution} : J'ai dédramatisé leur rôle en insistant sur le fait qu'ils étaient des "témoins" et non des "juges", et que je portais la responsabilité finale de la note.
    
    \item \textbf{Pression du groupe} : Tentatives d'intimidation ou de "marchandage" pour obtenir une répartition égalitaire.
    \item \textbf{Solution} : Sanctuarisation de la "Table Ronde" comme un espace protégé où la parole du chef est libre et soutenue par l'enseignant.
    
    \item \textbf{Biais de sévérité} : Paradoxalement, certains chefs étaient plus durs avec leurs amis pour prouver leur impartialité.
    \item \textbf{Solution} : Mon rôle de régulateur a consisté parfois à... remonter des notes proposées trop basses !
\end{itemize}

\subsubsection{Axes d'amélioration}
Pour perfectionner ce dispositif \parencite{ingenierie_digiforma}, je prévois :
\begin{itemize}
    \item \textbf{Formation explicite au feedback} : Consacrer une séance à apprendre "comment dire les choses qui fâchent" de manière constructive (Communication Non Violente).
    \item \textbf{Feedback à 360°} : Permettre aux membres d'évaluer aussi leur chef de projet, pour équilibrer le pouvoir.
    \item \textbf{Anonymisation partielle} : Tester un système où les propositions de notes sont d'abord faites anonymement avant d'être discutées.
\end{itemize}

\subsection{Transférabilité et validation des compétences}

Cette expérience apporte des preuves probantes pour la validation conjointe des UEs/ECs du Master MEEF PIF (Parcours FFMS d'Unicaen) et des Blocs de Compétences nationaux (RNCP38152) :

\subsubsection{EC 151 / 251 (Analyse de l'activité) | RNCP BC06 (Praticien réflexif)}

Conception et expérimentation de modalités évaluatives novatrices répondant à des enjeux didactiques complexes \parencite{schon_reflective_practitioner,perrenoud_pratique_reflexive,recherche_action_education}.

\subsubsection{EC 162 / 262 (Numérique éducatif) | RNCP BC05 (Évaluation formative)}

Utilisation de l'évaluation comme un \textbf{levier d'apprentissage} pour développer l'autonomie, la responsabilité et les compétences sociales des apprenants \parencite{evaluation_formative}.

\subsubsection{EC 125 / 225 (Évaluation des projets) | RNCP BC04 (Ingénierie de formation \& IUT)}

Cette méthode d'auto-évaluation et de régulation s'inspire des pratiques managériales professionnelles \parencite{competences_transversales_eurecia}, particulièrement adaptées aux étudiants de BUT et à la formation d'adultes.
